%% kapa: exam=1 round=30 year=2019 period=1 form=A subject=civil_law subject_ko=민법 seq=4 number=4
%% source: Q-Net 감정평가사 기출문제 (cst00309) artlSeq=5205205 file=extracted/1008/5205205/2019년 제30회 감정평가사 1차시험 문제지/2019년 제30회 감정평가사 1차시험 문제지/2019년도 제30회 감정평가사 1차 1교시 A형.hwp
%% conversion: kordoc 4.24.51 -> tools/gen1.py (auto); review=proofread-sonnet; flags=none
\begin{question}{4}
18세인 미성년자가 단독으로 유효하게 할 수 있는 행위가 아닌 것은?
\begin{choices}
  \item 자신이 제한행위능력자임을 이유로 취소할 수 있는 법률행위의 취소
  \item 부모로부터 받은 한 달분의 용돈을 친구에게 빌려주는 행위
  \item 자전거를 부담부로 증여받는 행위
  \item 타인의 대리인으로서 토지를 매도하는 행위
  \item 부모의 동의를 받아 법률상 혼인을 한 후, 주택을 구입하는 행위
\end{choices}
\end{question}
\answer{3}
